\documentclass{amsart}
\usepackage{amssymb, color}
\usepackage{mathrsfs}
\usepackage{esint}
\usepackage{enumerate}

\usepackage[colorlinks, citecolor=blue,pagebackref,hypertexnames=false]{hyperref}


\numberwithin{equation}{section}

\begin{document}


\newtheorem{cl}{Claim}
\newtheorem{theorem}{Theorem}[section]
\newtheorem{proposition}[theorem]{Proposition}
\newtheorem{coro}[theorem]{Corollary}
\newtheorem{lemma}[theorem]{Lemma}
\newtheorem{definition}[theorem]{Definition}
\newtheorem{nota}[theorem]{Notation}
\newtheorem{assum}{Assumption}[section]
\newtheorem{example}[theorem]{Example}
\newtheorem{remark}[theorem]{Remark}
\newtheorem{construction}[theorem]{Construction}


\title[Schatten class properties of Dunkl--Riesz transform commutators]{Schatten class properties of Dunkl--Riesz transform commutators}



\author{Zhijie Fan}
\address{Zhijie Fan, School of Mathematical Sciences, South China Normal University, Guangzhou 510631, China}
\email{fanzhj3@mail2.sysu.edu.cn}

\author{Ji Li}
\address{Ji Li, Department of Mathematics, Macquarie University, NSW 2109, Australia}
\email{ji.li@mq.edu.au}


\author{Dmitriy Zanin}
\address{School of Mathematics and Statistics, Central South University, Changsha 410075, People's Republic of China}
\email{d.zanin@csu.edu.cn}






  \date{\today}

 \subjclass[2010]{47B10, 42B20, 43A85}
\keywords{Schatten class, Riesz transform commutators, Dunkl operator, Besov space, Nearly weakly orthogonal sequence}




\begin{abstract}

\end{abstract}

\maketitle


\tableofcontents

\section{Introduction}
\subsection{Background and motivation}



\begin{theorem}\label{main theorem} Let $R$ be a normalised root system in $\mathbb{R}^N$ and let $G$ be the corresponding finite reflection group. Let $k:R\to\mathbb{R}_+$ be $G$-invariant multiplicity function. Let $R_{{\rm Dunkl},j}$ be the $j$-th Dunkl-Riesz transform corresponding to $R$ and $k.$ Let $w$ be the Dunkl weight corresponding to $R$ and $k.$ Let $f\in L_1^{{\rm loc}}(\mathbb{R}^N).$
\begin{enumerate}[{\rm (i)}]
\item\label{mta} If $f\in\dot{W}^{\frac{N}{p},p}(\mathbb{R}^N),$ $p>N,$ then $[R_{{\rm Dunkl},j},M_f]\in\mathcal{L}_p(L_2(\mathbb{R}^N,w(x)dx)).$ Furthermore,
$$\|[R_{{\rm Dunkl},j},M_f]\|_{\mathcal{L}_p(L_2(\mathbb{R}^N,w(x)dx))}\leq c_{p,R,k}\|f\|_{\dot{W}^{\frac{N}{p},p}(\mathbb{R}^N)};$$
\item\label{mtb} If $[R_{{\rm Dunkl},j},M_f]\in\mathcal{L}_p(L_2(\mathbb{R}^N,w(x)dx)),$ $p>N,$ then $f\in\dot{W}^{\frac{N}{p},p}(\mathbb{R}^N).$ Furthermore,
$$\|[R_{{\rm Dunkl},j},M_f]\|_{\mathcal{L}_p(L_2(\mathbb{R}^N,w(x)dx))}\leq c_{p,R,k}^{-1}\|f\|_{\dot{W}^{\frac{N}{p},p}(\mathbb{R}^N)};$$
\item\label{mtc} If $f\in\dot{W}^{1,N}(\mathbb{R}^N),$ then $[R_{{\rm Dunkl},j},M_f]\in\mathcal{L}_{N,\infty}(L_2(\mathbb{R}^N,w(x)dx)).$ Furthermore,
$$\|[R_{{\rm Dunkl},j},M_f]\|_{\mathcal{L}_{N,\infty}(L_2(\mathbb{R}^N,w(x)dx))}\leq c_{R,k}\|f\|_{\dot{W}^{1,N}(\mathbb{R}^N)};$$
\item\label{mtd} If $[R_{{\rm Dunkl},j},M_f]\in\mathcal{L}_{N,\infty}(L_2(\mathbb{R}^N,w(x)dx)),$ then $f\in\dot{W}^{1,N}(\mathbb{R}^N).$ Furthermore,
$$\|[R_{{\rm Dunkl},j},M_f]\|_{\mathcal{L}_{N,\infty}(L_2(\mathbb{R}^N,w(x)dx))}\geq c_{R,k}^{-1}\|f\|_{\dot{W}^{1,N}(\mathbb{R}^N)};$$
\end{enumerate}
\end{theorem}

\subsection{Strategy of the proof}

In Section \ref{comparison section}, we introduce our {\it ad hoc} Sobolev space and demonstrate that its semi-norm is controlled by the usual Sobolev semi-norm. In Section \ref{upper estimate section}, we use Janson-Wolff method to control the $\|\cdot\|_{\mathcal{L}_p(L_2(\mathbb{R}^N,w(x)dx))}$-norm of the commutator $[R_{{\rm Dunkl},j},M_f]$ by the {\it ad hoc} Sobolev semi-norm of $f.$ Thus, we control the $\|\cdot\|_{\mathcal{L}_p(L_2(\mathbb{R}^N,w(x)dx))}$-norm of the commutator $[R_{{\rm Dunkl},j},M_f]$ by the usual Sobolev semi-norm (i.e., we prove Theorem \ref{main theorem} \eqref{mta}).

In Section \ref{lower estimate section}, we control the oscillation semi-norm of $f$ (with respect to measure $w(x)dx$ and distance $(x,y)\to|x-y|_{\infty}$) via $\|\cdot\|_{\mathcal{L}_p(L_2(\mathbb{R}^N,w(x)dx))}$-norm of the commutator $[R_{{\rm Dunkl},j},M_f].$ After that, we verify the assumptions in Proposition 1.2 in \cite{Hytonen-new} for measures $w(x)dx$ and $dx$ and, using the latter proposition, conclude that oscillation semi-norm of $f$ with respect to measure $w(x)dx$ and distance $(x,y)\to|x-y|_{\infty}$ is equivalent to the one with respect to measure $dx$ and distance $(x,y)\to|x-y|_{\infty}.$ The latter is known to be equivalent to the usual Sobolev semi-norm of $f.$

\subsection{Notation}
Conventionally, we let $\mathbb{N}$ denote the set of positive integers and set $\mathbb{Z}_+ := \{0\} \cup \mathbb{N}$. Throughout the paper, the symbols $C$ and $c$ denote positive constants (possibly varying from line to line) that may depend only on $N$, $p$, and the structural data $(R,k)$ of the Dunkl setting, but are independent of the main parameters (e.g., the symbol $b$ and auxiliary decay factors).
 We write $A \lesssim B$ to mean that $A \leq C B$ for some constant $C > 0$, and $A \sim B$ to mean that both $A \lesssim B$ and $B \lesssim A$ hold.
For \(\lambda>0\) and a cube \(Q\) with sidelength \(\ell(Q)\), write \(\lambda Q\) for the cube concentric with \(Q\) and sidelength \(\lambda\,\ell(Q)\). Similarly, for every ball \(B=B(x,r)\), set \(\lambda B:=B(x,\lambda r)\). For any $1\leq p\leq\infty$, we denote by $p'$ the conjugate of $p$, which satisfies $\frac{1}{p}+\frac{1}{p'}=1$. We denote the average of a function $f$ over a $\mu$-measurable set $E$ by
\begin{align}\label{defaverage}
(f)_E:=\fint_{E}f(x)d\mu(x):=\frac{1}{\mu(E)}\int_{E}f(x)d\mu(x).
\end{align}
For a subset $E \subset \mathbb{R}^{N}$, we denote by $\chi_E$ its characteristic function.

Denote by $\{e_j\}_{j=1}^N$ the standard orthonormal basis on $\mathbb{R}^N$.
%For two measurable Borel sets $E$ and $F$, we use the notation ${\rm dist}(E,\,F)$ to denote the distance between them.


\section{Preliminaries}
\setcounter{equation}{0}
\subsection{Harmonic analysis in the Dunkl setting}
Let $\mathbb{R}^N$ be the Euclidean space equipped with the usual scalar product $\langle \cdot,\,\cdot\rangle $ and the induced norm $\|\cdot\|$. For a nonzero vector $v\in \mathbb{R}^N$, the reflection $\sigma_v$ with respect to the hyperplane orthogonal to a normalised vector $v$ is defined by
\begin{align*}
	\sigma_v x:= x-2\langle x,v\rangle v.
\end{align*}
A finite set $R\subset \mathbb{R}^N\backslash \{0\}$ is called a root system if $\sigma_v(R)=R$ for every $v\in R$. A root system $R$ is called a normalized reduced root system if  $|v|=1$ for every $v\in R.$
The finite group $G$ generated by the reflections $\{\sigma_v\}_{v\in R}$ is called the Weyl group (reflection group) of the root system $R$.
Let $k: R \to [0,\infty)$ be a multiplicity function, which is a $G$-invariant function  fixed throughout the paper.

Define the Dunkl measure $w(x)dx$ on $\mathbb{R}^N$ by
$$w(x)=\prod_{v \in R} |\langle x,v\rangle|^{k(v)}.$$
Let $B(x,r):=\{y\in \mathbb{R}^N:\ |x-y|< r\}$ be the Euclidean ball centred at $x$ with radius $r>0$.
Furthermore,
\begin{align}\label{twosideinequ}
  w(B(x,r))\sim r^N \prod_{v \in R}(|\langle x, v \rangle |+r)^{k(v)},
\end{align}
where the implicit constant only depends on $N,R$ and $k$.  It follows from \eqref{twosideinequ} that there exists a constant $C\geq 1$ such that for every $x\in \mathbb{R}^N$,
\begin{align}\label{doublinginequality}
C^{-1}\left(\frac{r_2}{r_1}\right)^N\leq \frac{w(B(x,r_2))}{w(B(x,r_1))}\leq C\left(\frac{r_2}{r_1}\right)^{N'},\ {\rm for}\ 0<r_1<r_2,
\end{align}
so the numbers $N'$ and $N$ are called the upper and lower homogeneous dimensions, respectively.

The Dunkl operators $T_\xi$ are defined by
$$T_\xi f(x):=\partial_\xi f(x)+\sum_{v\in R}\frac{k(v)}{2}\langle v,\xi\rangle\frac{f(x)-f(\sigma_v(x))}{\langle v,x\rangle},$$
where $\partial_\xi f$ denotes the directional derivative of $f$ in the direction $\xi$. Furthermore, the Dunkl Laplacian is defined by
$$\Delta_{{\rm Dunkl}}=-\sum_{j=1}^N T_{e_j}^2.$$
In this article, we consider the Dunkl-Riesz transforms defined by $R_{{\rm Dunkl},j}:=T_{e_j}\Delta_{{\rm Dunkl}}^{-\frac12}$, $j=1,2,\cdots,N.$


\begin{definition}{\rm
The Weyl chambers are defined as the closures of connected components of the set
  \begin{align}
    \left\{x\in\mathbb{R}^N:\langle x,v\rangle \neq 0\ \ { for} \ { all} \ v\in R\right\}.\end{align}}
\end{definition}
Define
$$d(x,y)=\min_{\sigma\in G}|x-\sigma(y)|,\quad x,y\in\mathbb{R}^N.$$
Recall from \cite[Chapter VII, proof of Theorem 2.12]{MR514561} that
\begin{align}\label{distanceequa}
d(x,y)=|x-y|\mbox{ if and only if }x,y\in\Gamma\mbox{ for some chamber }\Gamma.
\end{align}



Denote by $K_{{\rm Dunkl},j}(x,y)$ the integral kernel of $R_{{\rm Dunkl},j}$. We recall the kernel estimates for the Riesz transforms in the Dunkl setting proved in \cite{MR4733961}. Although these estimates may not be sharp (see the comment following Theorem~2.1 in \cite[p.~10]{MR4733961}), they suffice for our purposes.
\begin{lemma}\label{kernelestimate}
There exists a constant $C$ such that for $j\in\{1,2,\cdots,N\}$ and every $x,y$ with $d(x,y)\neq 0,$
\begin{align*}
|K_{{\rm Dunkl},j}(x,y)|\leq \frac{d(x,y)}{|x-y|}\cdot\frac{C}{w(B(x,d(x,y)))}.
\end{align*}
\end{lemma}
\begin{proof}
The proof is given in \cite[Theorem 1.1]{MR4733961}.
\end{proof}

\begin{lemma}\label{for sufficiency auxiliary lemma} We have
$$d(x,y)^2\cdot w(x)w(y)\leq Cw^2(B(x,d(x,y))),\quad x,y\in\mathbb{R}^N.$$
\end{lemma}
\begin{proof} Define equivalence relation $\sim$ on $R$ by saying $v_1\sim v_2$ if there exists $g\in G$ with $v_1=gv_2.$ Let us split $R$ into equivalence classes with respect to $\sim.$ Since $k$ is $G$-invariant, it follows that $k$ is constant on each equivalence class. Note that $G$ leaves each equivalence class invariant.

We, therefore, write
$$w(x)=\prod_{A\in R/\sim}\big(\prod_{v\in A}|\langle x,v\rangle|\big)^{k|_A}.$$

On the other hand,
$$w(B(x,r))\geq c_{R,k}^{(1)}r^N\prod_{A\in R/\sim}\big(\prod_{v\in A}(r+|\langle x,v\rangle|)\big)^{k|_A}.$$

Let $x\in\Gamma$ and $y\in g\Gamma.$ Set $z=g^{-1}y\in\Gamma.$ Note that $d(x,y)=|x-z|$ and, therefore,
$$|\langle x,v\rangle|+d(x,y)=|\langle x,v\rangle|+|x-z|\geq 2^{-\frac12}|\langle z,v\rangle|.$$
Thus,
$$\prod_{v\in A}(|\langle x,v\rangle|+d(x,y))\geq\prod_{v\in A}2^{-\frac12}|\langle z,v\rangle|=2^{-\frac12|A|}\prod_{v\in A}|\langle g^{-1}y,v\rangle|=$$
$$=2^{-\frac12|A|}\prod_{v\in A}|\langle y,gv\rangle|=2^{-\frac12|A|}\prod_{v\in gA}|\langle y,v\rangle|=2^{-\frac12|A|}\prod_{v\in A}|\langle y,v\rangle|.$$
Thus,
$$w(B(x,d(x,y)))\geq c_{R,k}^{(1)}d^N(x,y)\prod_{A\in R/\sim}\big(\prod_{v\in A}(d(x,y)+|\langle x,v\rangle|)\big)^{k|_A}\geq$$
$$\geq c_{R,k}^{(1)}d^N(x,y)\prod_{A\in R/\sim}\big(2^{-\frac12|A|}\prod_{v\in A}|\langle y,v\rangle|\big)^{k|_A}=c_{R,k}^{(2)}d^N(x,y)w(y).$$
On the other hand,
$$w(B(x,d(x,y)))\geq c_{R,k}^{(1)}d^N(x,y)w(x).$$
Multiplying the last two equations, we complete the proof.
\end{proof}

\begin{lemma}\label{nondegenerate} Let $1\leq j\leq N.$ If $x_j\geq y_j,$ then
$$K_{{\rm Dunkl},j}(x,y)\geq c_{R,k}\frac{x_j-y_j}{|x-y|}\cdot\frac1{w(B(x,|x-y|))}.$$
If $x_j\leq y_j,$ then
$$K_{{\rm Dunkl},j}(x,y)\leq c_{R,k}\frac{x_j-y_j}{|x-y|}\cdot\frac1{w(B(x,|x-y|))}.$$
\end{lemma}
\begin{proof} The lower bound is identical to that in \cite[Theorem 1.2]{MR4733961}.
\end{proof}

\subsection{Schatten class}
{\color{red}[Zhijie will add this later]}





\subsection{Nearly weakly  orthonormal sequences}
Rochberg and Semmes \cite{RS} introduced the notion of \emph{nearly weakly  orthonormal (NWO)} sequences of functions. This notion has since proved highly effective in characterizing Schatten class membership of  singular integral commutators in a variety of contexts (see \cite{FLLXarxiv,HKarxiv,MR4756021,MR4552581,MR4873796,LLW,RS,Zhangwei1,Zhangwei2}). The properties of NWO sequences were originally investigated in the Euclidean setting (equipped with Lebesgue measure), but the extension to spaces of homogeneous type is straightforward once a system of dyadic cubes is available in this general framework (see \cite[Section 13]{HKarxiv} for a complete proof).  For our purposes, we adapt  the extension to the Dunkl setting $(\mathbb{R}^N,\|\cdot\|,w)$ rather than recalling the general  definition
of NWO sequences and the general extension to spaces of homogeneous type.

\begin{definition} Let $\mathcal{D}$ be the dyadic lattice in $\mathbb{R}^N.$ A sequence of functions $\{e_Q\}_{Q\in\mathcal{D}}$ is called nearly weakly orthonormal (NWO, in short) sequence if
\begin{enumerate}[{\rm (i)}]
\item ${\rm supp}(e_Q)\subset cQ$ for some positive constant $c$ independent of $Q.$
\item $\|e_Q\|_{L_{\infty}(\mathbb{R}^N,w(x)dx)}\leq w(Q)^{-\frac12}.$
\end{enumerate}
\end{definition}

\begin{lemma}\label{NWO lower estimate}
Let $1<p<\infty$ and let $\mathcal{D}$ be the dyadic lattice in $\mathbb{R}^N.$
If $T\in\mathcal{L}_p(L_2(\mathbb{R}^N,w(x)dx)),$ then
$$\Big\|\Big\{\langle T e_Q,f_Q \rangle_{L_2(\mathbb{R}^N,w(x)dx)}\Big\}_{Q\in\mathcal{D}}\Big\|_{l_{p,q}(\mathcal{D})}\leq c_{N,p,q,w,c_1,c_2}\|T\|_{\mathcal{L}_{p,q}(L_2(\mathbb{R}^N,w(x)dx))}$$
for every NWO systems $\{e_Q\}_{Q\in\mathcal{D}}$ and $\{f_Q\}_{Q\in\mathcal{D}}$ (with respective constants $c_1$ and $c_2$).
\end{lemma}
\begin{proof} {\color{red} reference to \cite{RS} is inappropriate. Rochberg and Semmes do not consider weighted setting. we absolutely must provide correct reference.} This property can be found in  \cite[(1.10)]{RS} (see also \cite[Section 7]{HKarxiv}).
\end{proof}


\section{Comparison of Sobolev norms}\label{comparison section}

\begin{theorem}\label{sobolev comparison theorem} For every $f\in\dot{W}^{\frac{N}{p},p}(\mathbb{R}^N),$ $p>N,$ we have
$$\left(\int_{\mathbb{R}^N\times\mathbb{R}^N}\Big(\frac{d(x,y)}{|x-y|}\Big)^p|f(x)-f(y)|^p \frac{dxdy}{d(x,y)^{2N}}\right)^{\frac1p}\leq c_{p,R}\|f\|_{\dot{W}^{\frac{N}{p},p}(\mathbb{R}^N)}.$$
\end{theorem}

\begin{lemma}\label{s4 first compute lemma} For every convex polyhedral cone $\Gamma,$ we have
$$\int_{\Gamma}\frac{dx}{|x-y|^{2N}}\geq c_{\Gamma}{\rm dist}(y,\Gamma)^{-N},\quad y\notin\Gamma.$$
\end{lemma}
\begin{proof} Find a point $z\in\Gamma$ such that
$$|y-z|={\rm dist}(y,\Gamma).$$
We have $z+\Gamma\subset\Gamma.$ Thus,
$$\int_{\Gamma}\frac{dx}{|x-y|^{2N}}\geq \int_{z+\Gamma}\frac{dx}{|x-y|^{2N}}=\int_{\Gamma}\frac{dx}{|x-(y-z)|^{2N}}.$$
Clearly,
$$|x-(y-z)|\leq |x|+|y-z|.$$
Thus,
$$\int_{\Gamma}\frac{dx}{|x-y|^{2N}}\geq\int_{\Gamma}\frac{dx}{(|x|+|y-z|)^{2N}}=|y-z|^{-N}\int_{\Gamma}\frac{dx}{(|x|+1)^{2N}}.$$
\end{proof}

\begin{lemma}\label{s4 second compute lemma} If $p\in(N,2N),$ then
$$\sup_{a>0}\int_{\mathbb{R}^N}\frac{|x-ae_N|^{p-2N}dx}{(1+|x|)^p}<\infty.$$
\end{lemma}
\begin{proof} {\bf Step 1:} we claim that, for every $a>0,$
$$\int_{\mathbb{R}^N}\frac{|x-ae_N|^{p-2N}dx}{(1+|x|)^p}\leq$$
$$\leq 4^p\int_{\mathbb{R}^{N-1}\times(-a,a)}\frac{|x|^{p-2N}dx}{(1+|x'|+a)^p}+3^{2N-p}\int_{\mathbb{R}^N}\frac{|x|^{p-2N}dx}{(1+|x|)^p}.$$


We clearly have
$$\int_{\mathbb{R}^N}\frac{|x-ae_N|^{p-2N}dx}{(1+|x|)^p}=$$
$$=\int_{\mathbb{R}^{N-1}\times(\frac{a}{2},\frac{3a}{2})}\frac{|x-ae_N|^{p-2N}dx}{(1+|x|)^p}+\int_{\mathbb{R}^{N-1}\times\mathbb{R}\backslash(\frac{a}{2},\frac{3a}{2})}\frac{|x-ae_N|^{p-2N}dx}{(1+|x|)^p}.$$

If $x_N\notin(\frac{a}{2},\frac{3a}{2}),$ then $|x_N-a|\geq \frac13|x_N|.$ Hence, $|x-ae_N|\geq \frac13|x|.$ Thus,
$$\int_{\mathbb{R}^{N-1}\times\mathbb{R}\backslash(\frac{a}{2},\frac{3a}{2})}\frac{|x-ae_N|^{p-2N}dx}{(1+|x|)^p}\leq$$
$$\leq \int_{\mathbb{R}^{N-1}\times\mathbb{R}\backslash(\frac{a}{2},\frac{3a}{2})}\frac{3^{2N-p}|x|^{p-2N}dx}{(1+|x|)^p}\leq \int_{\mathbb{R}^N}\frac{3^{2N-p}|x|^{p-2N}dx}{(1+|x|)^p}.$$

If $x_N\in(\frac{a}{2},\frac{3a}{2}),$ then (here, we using the notation $x'=(x_1,\cdots,x_{N-1})$) $$1+|x|\geq1+\max\{|x'|,\frac12a\}\geq\frac14(1+|x'|+a).$$
Thus,
$$\int_{\mathbb{R}^{N-1}\times(\frac{a}{2},\frac{3a}{2})}\frac{|x-ae_N|^{p-2N}dx}{(1+|x|)^p}\leq 4^p\int_{\mathbb{R}^{N-1}\times(\frac{a}{2},\frac{3a}{2})}\frac{|x-ae_N|^{p-2N}dx}{(1+|x'|+a)^p}=$$
$$=4^p\int_{\mathbb{R}^{N-1}\times(-\frac{a}{2},\frac{a}{2})}\frac{|x|^{p-2N}dx}{(1+|x'|+a)^p}\leq 4^p\int_{\mathbb{R}^{N-1}\times(-a,a)}\frac{|x|^{p-2N}dx}{(1+|x'|+a)^p}.$$

Combining all preceding estimates, we infer the claim in Step 1.

{\bf Step 2:} we claim that
$$\sup_{a>0}\int_{\mathbb{R}^{N-1}\times(-a,a)}\frac{|x|^{p-2N}dx}{(1+|x'|+a)^p}<\infty.$$

Passing to polar coordinates in $\mathbb{R}^{N-1},$ we equivalently rewrite the claim as
$$\sup_{a>0}\int_{\mathbb{R}_+\times(0,a)}\frac{(r^2+x_N^2)^{\frac{p}{2}-N}r^{N-2}drdx_N}{(1+r+a)^p}<\infty.$$

If $0<a<1,$ then
$$\int_{\mathbb{R}_+\times(0,a)}\frac{(r^2+x_N^2)^{\frac{p}{2}-N}r^{N-2}drdx_N}{(1+r+a)^p}\leq \int_{\mathbb{R}_+\times(0,1)}\frac{(r^2+x_N^2)^{\frac{p}{2}-N}r^{N-2}drdx_N}{(1+r)^p}.$$

If $a\geq 1,$ then
$$\int_{\mathbb{R}_+\times(0,a)}\frac{(r^2+x_N^2)^{\frac{p}{2}-N}r^{N-2}drdx_N}{(1+r+a)^p}\leq \int_{\mathbb{R}_+\times(0,a)}\frac{(r^2+x_N^2)^{\frac{p}{2}-N}r^{N-2}drdx_N}{(r+a)^p}=$$
$$=\int_{(a,\infty)\times(0,a)}\frac{(r^2+x_N^2)^{\frac{p}{2}-N}r^{N-2}drdx_N}{(r+a)^p}+$$
$$+\int_0^a\int_0^r\frac{(r^2+x_N^2)^{\frac{p}{2}-N}r^{N-2}dx_Ndr}{(r+a)^p}+\int_0^a\int_0^{x_N}\frac{(r^2+x_N^2)^{\frac{p}{2}-N}r^{N-2}drdx_N}{(r+a)^p}\leq$$
$$\leq\footnote{It is exactly at this point we are using the assumption $p<2N.$}\int_{(a,\infty)\times(0,a)}\frac{r^{p-2N}r^{N-2}drdx_N}{r^p}+$$
$$+\int_0^a\int_0^r\frac{r^{p-2N}r^{N-2}dx_Ndr}{a^p}+\int_0^a\int_0^{x_N}\frac{x_N^{p-2N}r^{N-2}drdx_N}{a^p}=$$
$$=\int_a^{\infty}r^{-N-1}dr+a^{-p}\int_0^ar^{p-N-1}dr+\frac1{N-1}a^{-p}\int_0^ax_N^{p-N-1}dx_N=$$
$$=\frac1Na^{-N}+\frac1{p-N}a^{-N}+\frac1{(N-1)(p-N)}a^{-N}\leq \frac1N+\frac1{p-N}+\frac1{(N-1)(p-N)}.$$
This yields the claim in Step 2.

Combining Steps 1 and 2, we complete the proof of the lemma.
\end{proof}


\begin{lemma}\label{s4 third compute lemma} If $p\in(N,2N),$ then $$\int_{\mathbb{R}^{N-1}\times\mathbb{R}_+}\frac{|x-y|^{p-2N}dx}{|x+z|^p}\leq c_{p,N}z_N^{-N},\quad y,z\in\mathbb{R}^{N-1}\times\mathbb{R}_+.$$
\end{lemma}
\begin{proof} By dilation invariance, we may assume without loss of generality that $z_N=1.$ By translation invariance in $\mathbb{R}^{N-1},$ we may assume without loss of generality that $z_1=\cdots=z_{N-1}=0.$ We, therefore, have
$$|x+z|=|x+e_N|\geq \max\{1,|x|\}\geq\frac12(1+|x|),\quad x\in\mathbb{R}^{N-1}\times\mathbb{R}_+.$$
Thus,
$$\int_{\mathbb{R}^{N-1}\times\mathbb{R}_+}\frac{|x-y|^{p-2N}dx}{|x+z|^p}\leq 2^p\int_{\mathbb{R}^{N-1}\times\mathbb{R}_+}\frac{|x-y|^{p-2N}dx}{(1+|x|)^p}\leq$$
$$=2^p\int_{\mathbb{R}^N}\frac{|x-y|^{p-2N}dx}{(1+|x|)^p}=2^p\int_{\mathbb{R}^N}\frac{\big|x-|y|e_N\big|^{p-2N}dx}{(1+|x|)^p}.$$
The assertion follows from Lemma \ref{s4 second compute lemma}.
\end{proof}


We need the following elementary version of the Hahn-Banach theorem.

\begin{lemma}\label{hahn-banach lemma} Let $K\subset\mathbb{R}^N$ be a convex set and let $y\notin K.$ Let $z\in K$ be the point closest to $y.$ Affine hyperplane $P$ passing through $z$ and orthogonal to $y-z$ splits the space into two half-spaces, one of them containing $y,$ the other one containing $K.$
\end{lemma}
\begin{proof} By translation invariance, we may assume without loss of generality that $z=0.$ By rotation and dilation invariance, we may assume without loss of generality that $y=e_N.$ By assumption, ${\rm dist}(e_N,K)=1.$ We need to show that $K$ is contained in the (closed) lower half-plane.

Assume the contrary: let $x\in K$ be such that $x_N>0.$ Since $0\in K$ and since $K$ is convex, it follows that $\epsilon x\in K$ for every $\epsilon\in(0,1).$ Thus,
$${\rm dist}(e_N,K)\leq |e_N-\epsilon x|,\quad \epsilon\in(0,1).$$
Since $\epsilon\in(0,1)$ is arbitrary, it follows that
$${\rm dist}(e_N,K)\leq \inf_{\epsilon\in(0,1)}|e_N-\epsilon x|=\inf\big(\epsilon^2\sum_{k=1}^{N-1}x_k^2+(1-\epsilon x_N)^2\big)^{\frac12}.$$
Since $x_N>0,$ it follows that infimum is strictly less than $1.$ Thus, ${\rm dist}(e_N,K)<1.$ However, it is granted in the lemma that ${\rm dist}(e_N,K)=1.$ This contradiction shows that such an $x$ cannot exist. This completes the proof.
\end{proof}


\begin{lemma}\label{s4 main lemma} If $p\in(N,2N),$ then
$$\int_{\Gamma}\frac{|x-y|^{p-2N}dx}{|x-gy|^p}\leq c_p{\rm dist}(gy,\Gamma)^{-N},\quad y\in\Gamma.$$
\end{lemma}
\begin{proof} Find a point $z\in\Gamma$ such that
$$|gy-z|={\rm dist}(gy,\Gamma).$$
By Lemma \ref{hahn-banach lemma}, affine hyperplane passing through $z$ and orthogonal to $gy-z$ splits the space into two half-spaces, one of them containing $gy,$ the other one containing $\Gamma$ (denote this one by $H$).

We have $y\in H,$ $\Gamma\subset H$ and $gy\notin H.$ So,
$$\int_{\Gamma}\frac{|x-y|^{p-2N}dx}{|x-gy|^p}\leq\int_H\frac{|x-y|^{p-2N}dx}{|x-gy|^p}.$$
Using Lemma \ref{s4 third compute lemma}, we write
$$\int_H\frac{|x-y|^{p-2N}dx}{|x-gy|^p}\leq c_{p,N}{\rm dist}(gy,H)^{-N}.$$
Since $z\in \partial H$ and since $gy-z$ is orthogonal to $\partial H,$ it follows that ${\rm dist}(gy,H)=|gy-z|={\rm dist}(gy,\Gamma).$
\end{proof}


\begin{proof}[Proof of Theorem \ref{sobolev comparison theorem}] {\bf Case 1:} Let $p\in(N,2N).$

We have
$$\Big(\int_{\Gamma\times g\Gamma}\Big(\frac{d(x,y)}{|x-y|}\Big)^p|f(x)-f(y)|^p \frac{dxdy}{d(x,y)^{2N}}\Big)^{\frac1p}=$$
$$=\Big(\int_{\Gamma\times \Gamma}\Big(\frac{|x-y|}{|x-gy|}\Big)^p|f(x)-f(gy)|^p \frac{dxdy}{|x-y|^{2N}}\Big)^{\frac1p}\leq$$
$$\leq\Big(\int_{\Gamma\times \Gamma}\Big(\frac{|x-y|}{|x-gy|}\Big)^p|f(x)-f(y)|^p \frac{dxdy}{|x-y|^{2N}}\Big)^{\frac1p}+$$
$$+\Big(\int_{\Gamma\times \Gamma}\Big(\frac{|x-y|}{|x-gy|}\Big)^p|f(y)-f(gy)|^p \frac{dxdy}{|x-y|^{2N}}\Big)^{\frac1p}\leq$$
$$\leq\Big(\int_{\Gamma\times \Gamma}|f(x)-f(y)|^p \frac{dxdy}{|x-y|^{2N}}\Big)^{\frac1p}+$$
$$+\Big(\int_{\Gamma}|f(y)-f(gy)|^p\Big(\int_{\Gamma}\frac{|x-y|^{p-2N}}{|x-gy|^p}dx\Big)dy\Big)^{\frac1p}.$$
Obviously, the first summand on the right hand side is less than $\|f\|_{\dot{W}^{\frac{N}{p},p}(\mathbb{R}^N)}.$ By Lemmas \ref{s4 main lemma} and \ref{s4 first compute lemma},
$$\Big(\int_{\Gamma\times g\Gamma}\Big(\frac{d(x,y)}{|x-y|}\Big)^p|f(x)-f(y)|^p \frac{dxdy}{d(x,y)^{2N}}\Big)^{\frac1p}\leq \|f\|_{\dot{W}^{\frac{N}{p},p}(\mathbb{R}^N)}+$$
$$+c_{p,R}^{(1)}\Big(\int_{\Gamma}|f(y)-f(gy)|^p\Big(\int_{\Gamma}\frac{dx}{|x-gy|^{2N}}\Big)dy\Big)^{\frac1p}=$$
$$=\|f\|_{\dot{W}^{\frac{N}{p},p}(\mathbb{R}^N)}+c_{p,R}^{\frac1p}\Big(\int_{\Gamma\times \Gamma}\frac{|f(y)-f(gy)|^pdxdy}{|x-gy|^{2N}}\Big)^{\frac1p}\leq$$
$$\leq\|f\|_{\dot{W}^{\frac{N}{p},p}(\mathbb{R}^N)}+c_{p,R}^{(1)}\Big(\int_{\Gamma\times \Gamma}\frac{|f(x)-f(gy)|^pdxdy}{|x-gy|^{2N}}\Big)^{\frac1p}+$$
$$+c_{p,R}^{(1)}\Big(\int_{\Gamma\times \Gamma}\frac{|f(x)-f(y)|^pdxdy}{|x-gy|^{2N}}\Big)^{\frac1p}.$$
Now,
$$\Big(\int_{\Gamma\times \Gamma}\frac{|f(x)-f(gy)|^pdxdy}{|x-gy|^{2N}}\Big)^{\frac1p}=\Big(\int_{\Gamma\times g\Gamma}\frac{|f(x)-f(y)|^pdxdy}{|x-y|^{2N}}\Big)^{\frac1p}\leq \|f\|_{\dot{W}^{\frac{N}{p},p}(\mathbb{R}^N)}$$
and
$$\Big(\int_{\Gamma\times \Gamma}\frac{|f(x)-f(y)|^pdxdy}{|x-gy|^{2N}}\Big)^{\frac1p}\leq \Big(\int_{\Gamma\times \Gamma}\frac{|f(x)-f(y)|^pdxdy}{|x-y|^{2N}}\Big)^{\frac1p}\leq \|f\|_{\dot{W}^{\frac{N}{p},p}(\mathbb{R}^N)}.$$
Thus,
$$\Big(\int_{\Gamma\times g\Gamma}\Big(\frac{d(x,y)}{|x-y|}\Big)^p|f(x)-f(y)|^p \frac{dxdy}{d(x,y)^{2N}}\Big)^{\frac1p}\leq (1+2c_{p,R}^{(1)}) \|f\|_{\dot{W}^{\frac{N}{p},p}(\mathbb{R}^N)}.$$
Summing over $\Gamma$ and $1\neq g\in G,$ we obtain
$$\left(\int_{\mathbb{R}^N\times\mathbb{R}^N}\Big(\frac{d(x,y)}{|x-y|}\Big)^p|f(x)-f(y)|^p \frac{dxdy}{d(x,y)^{2N}}\right)^{\frac1p}\leq |G|^2(1+2c_{p,R}^{(1)})\|f\|_{\dot{W}^{\frac{N}{p},p}(\mathbb{R}^N)}.$$

{\bf Case 2:} Let $p\geq 2N.$ We have
$$\Big(\int_{\Gamma\times g\Gamma}\Big(\frac{d(x,y)}{|x-y|}\Big)^p|f(x)-f(y)|^p \frac{dxdy}{d(x,y)^{2N}}\Big)^{\frac1p}\leq$$
$$\leq\Big(\int_{\Gamma\times g\Gamma}\Big(\frac{d(x,y)}{|x-y|}\Big)^{2N}|f(x)-f(y)|^p \frac{dxdy}{d(x,y)^{2N}}\Big)^{\frac1p}=\|f\|_{\dot{W}^{\frac{N}{p},p}(\mathbb{R}^N)}.$$
\end{proof}



\section{Proof of Theorem \ref{main theorem} \eqref{mta}}\label{upper estimate section}

Let $(X,\mu)$ be a measure space and let $L_p(L_{q,\infty})^{\rm symm}(X\times X,\mu\times\mu)$ denote the mixed-norm space consisting of all $\mu\times\mu$-measurable functions $f:X\times X \to \mathbb{C}$ such that
$$\int_X\|f(x,\cdot)\|_{L_{q,\infty}(X,\mu)}^pd\mu(x),\int_X\|f(\cdot,y)\|_{L_{q,\infty}(X,\mu)}^pd\mu(y)<\infty.$$
We equip $L_p(L_{q,\infty})^{\rm symm}(X\times X,\mu\times\mu)$ with its natural norm.

\begin{lemma}\label{Russo}
Let $(X,\mu)$ be a measure space and let $2<p<\infty.$ Suppose that $K\in L_p(L_{p',\infty})^{\rm symm}(X\times X,\mu\times\mu).$ The integral operator
$$T_Kf(x):=\int_X K(x,y)f(y)d\mu(y)$$
belongs to $S_{p,\infty}(L_2(X,d\mu)).$ Moreover, we have
$$\|T_K\|_{S_{p,\infty}(L_2(X,d\mu))}\leq c_p \|K\|_{L_p(L_{p',\infty})^{\rm symm}(X\times X,\mu\times\mu)}.$$
\end{lemma}
\begin{proof}
The proof is given in \cite[Lemma 2]{JW} (see also \cite[Corollary 11.2]{HKarxiv}).
\end{proof}

\begin{lemma}\label{s4 holder lemma} Let $p>2$ and let $\frac1q=1-\frac2p.$ We have $$\|F_1F_2\|_{L_p(L_{p',\infty})(X\times X,\mu\times\mu)}\leq 2^{\frac1{pp'}}\|F_1\|_{L_p(X\times X,\mu\times\mu)}\|F_2\|_{L_{\infty}(L_{q,\infty})(X\times X,\mu\times\mu)}.$$
\end{lemma}
\begin{proof} We clearly have
$$LHS=\Big(\int_X\|F_1(x,\cdot)F_2(x,\cdot)\|_{L_{p',\infty}(X,\mu)}^pd\mu(x)\Big)^{\frac1p}.$$
Clearly, $\frac1{p'}=\frac1p+\frac1q.$ Thus,
$$\|F_1(x,\cdot)F_2(x,\cdot)\|_{L_{p',\infty}(X,\mu)}
\leq 2^{\frac1{p'}}\|F_1(x,\cdot)\|_{L_{p,\infty}(X,\mu)}\|F_2(x,\cdot)\|_{L_{q,\infty}(X,\mu)}\leq$$
$$\leq 2^{\frac1{p'}}\|F_1(x,\cdot)\|_{L_p(X,\mu)}\|F_2(x,\cdot)\|_{L_{q,\infty}(X,\mu)}.$$
Hence,
$$LHS\leq 2^{\frac1{pp'}}\Big(\int_X\|F_1(x,\cdot)\|_{L_p(X,\mu)}^p\|F_2(x,\cdot)\|_{L_{q,\infty}(X,\mu)}^pd\mu(x)\Big)^{\frac1p}\leq$$
$$\leq 2^{\frac1{pp'}}\Big(\int_X\|F_1(x,\cdot)\|_{L_p(X,\mu)}^pd\mu(x)\Big)^{\frac1p}\times\sup_{x\in X}\|F_2(x,\cdot)\|_{L_{q,\infty}(X,\mu)}=$$
$$=2^{\frac1{pp'}}\|F_1\|_{L_p(X\times X,\mu\times\mu)}\|F_2\|_{L_{\infty}(L_{q,\infty})(X\times X,\mu\times\mu)}.$$
\end{proof}

\begin{lemma}\label{s4 first lemma}
Let $(X,\mu)$ be a measure space and let $2<p<\infty.$ Suppose that $F\in L_p(X\times X,H^2(\mu\times\mu))$ and $H\in L_{\infty}(L_{1,\infty})^{\rm symm}(X\times X,\mu\times\mu).$ We have
$$\|T_{FH}\|_{S_{p,\infty}(L_2(X,d\mu))}\leq c_p \|F\|_{L_p(X\times X,H^2(\mu\times\mu))}\|H\|_{L_{\infty}(L_{1,\infty})^{\rm symm}(X\times X,\mu\times\mu)}^{1-\frac2p}.$$
\end{lemma}
\begin{proof} By Lemma \ref{Russo},
$$\|T_{FH}\|_{S_{p,\infty}(L_2(X,d\mu))}\leq c_p \|FH\|_{L_p(L_{p',\infty})^{\rm symm}(X\times X,\mu\times\mu)}.$$

Let $\frac1q=1-\frac2p.$ It follows from Lemma \ref{s4 holder lemma} that
$$\|F_1F_2\|_{L_p(L_{p',\infty})^{{\rm symm}}(X\times X,\mu\times\mu)}\leq$$
$$\leq 2^{\frac1{pp'}}\|F_1\|_{L_p(X\times X,\mu\times\mu)}\|F_2\|_{L_{\infty}(L_{q,\infty})^{{\rm symm}}(X\times X,\mu\times\mu)}.$$
Setting $F_1=FH^{\frac2p}$ and $F_2=H^{\frac1q},$ we obtain
$$\|FH\|_{L_p(L_{p',\infty})^{\rm symm}(X\times X,\mu\times\mu)}\leq$$
$$\leq 2^{\frac1{pp'}}\|FH^{\frac2p}\|_{L_p(X\times X,\mu\times\mu)}\|H^{\frac1q}\|_{L_{\infty}(L_{q,\infty})^{{\rm symm}}(X\times X,\mu\times\mu)}=$$
$$=2^{\frac1{pp'}}\|F\|_{L_p(X\times X,H^2(\mu\times\mu))}\|H\|_{L_{\infty}(L_{1,\infty})^{{\rm symm}}(X\times X,\mu\times\mu)}^{\frac1q}=$$
$$=2^{\frac1{pp'}}\|F\|_{L_p(X\times X,H^2(\mu\times\mu))}\|H\|_{L_{\infty}(L_{1,\infty})^{{\rm symm}}(X\times X,\mu\times\mu)}^{1-\frac2p}.$$
\end{proof}

\begin{lemma}\label{s4 second lemma}
Let $(X,\mu)$ be a measure space and let $2<p<\infty.$ Suppose that $F\in L_p(X\times X,GH^2(\mu\times\mu))$ and $H\in L_{\infty}(L_{1,\infty})^{\rm symm}(X\times X,\mu\times\mu).$ We have
$$\|T_{FH}\|_{S_p(L_2(X,d\mu))}\leq c_{p,H}\|F\|_{L_p(X\times X,H^2(\mu\times\mu))}.$$
\end{lemma}
\begin{proof} The assertion follows from Lemma \ref{s4 first lemma} by real interpolation.
\end{proof}

The following proposition, strictly speaking, is not part of the proof of Theorem \ref{main theorem} \eqref{mta}. Instead, it features in the proof of Theorem \ref{main theorem} \eqref{mtb} in the next section. However, in terms of the employed technique, it certainly belongs to this section.

\begin{proposition}\label{upper bound}
Let $p>2,$ let $\Gamma$ be a chamber and let $f\in \dot{W}^{\frac{N}{p},p}(\Gamma).$ Let $V$ be an integral operator whose integral kernel $K$ satisfies the condition
$$|K(x,y)|\leq\frac1{w(B(x,|x-y|))}.$$
For every $1\leq j\leq N,$ we have
$$\|[V,M_f]\|_{\mathcal{L}_p(L_2(\Gamma,w(x)dx))}\leq c_{R,k}\|f\|_{\dot{W}^{\frac{N}{p},p}(\Gamma)}.$$
\end{proposition}
\begin{proof} In what follows, let $X=\Gamma$ and $d\mu(x)=w(x)dx.$ Set
$$F(x,y)=w(B(x,|x-y|))K(x,y)(f(x)-f(y)),\quad x\neq y,$$
$$H(x,y)=\frac1{w(B(x,|x-y|))},\quad x\neq y.$$
By \cite[Lemma 5.7]{HKarxiv}, $H\in L_{\infty}(L_{1,\infty})^{\rm symm}(\Gamma\times\Gamma,\mu\times\mu).$ By assumption on the kernel $K,$
$$\|F\|_{L_p(\Gamma\times\Gamma,G^2(\mu\times\mu))}\leq\Big(\int_{\Gamma\times\Gamma}|f(x)-f(y)|^p\frac{w(x)w(y)dxdy}{w^2(B(x,|x-y|))}\Big)^{\frac1p}.$$
By \eqref{twosideinequ},
$$w(B(x,|x-y|))\geq c_{R,k}'|x-y|^Nw^{\frac12}(x)w^{\frac12}(y),\quad x,y\in\mathbb{R}^N.$$
Thus,
$$\|F\|_{L_p(\Gamma\times\Gamma,H^2(\mu\times\mu))}\leq$$
$$\leq (c_{R,k}')^{-\frac2p}\Big(\int_{\Gamma\times\Gamma}\frac{|f(x)-f(y)|^p}{|x-y|^{2N}}dxdy\Big)^{\frac1p}=(c_{R,k}')^{-\frac2p}\|f\|_{\dot{W}^{\frac{N}{p},p}(\Gamma)}.$$
Obviously, the integral kernel of the operator $[V,M_f]$ is $FH.$ The assertion follows now from Lemma \ref{s4 second lemma}.
\end{proof}

\begin{proof}[Proof of Theorem \ref{main theorem} \eqref{mta}]
In what follows, let $X=\mathbb{R}^N$ and $d\mu(x)=w(x)dx.$ Set
$$F(x,y)=w(B(x,d(x,y)))K_{{\rm Dunkl},j}(x,y)(f(x)-f(y)),\quad x\neq y,$$
$$H(x,y)=\frac1{w(B(x,d(x,y)))},\quad x\neq y.$$
By \cite[Lemma 11.7]{HKarxiv}, $H\in L_{\infty}(L_{1,\infty})^{\rm symm}(X\times X,\mu\times\mu).$ By Lemma \ref{kernelestimate},
$$\|F\|_{L_p(X\times X,H^2(\mu\times\mu))}\leq$$
$$\leq c_{R,k}^{(1)}\left(\int_{\mathbb{R}^N\times\mathbb{R}^N}\Big(\frac{d(x,y)}{|x-y|}\Big)^p|f(x)-f(y)|^p\frac{w(x)w(y)dxdy}{w(B(x,d(x,y)))^2}\right)^{\frac1p}.$$
Using Lemma \ref{for sufficiency auxiliary lemma}, we conclude that
$$\|F\|_{L_p(X\times X,H^2(\mu\times\mu))}\leq c_{R,k}^{(2)}\left(\int_{\mathbb{R}^N\times\mathbb{R}^N}\Big(\frac{d(x,y)}{|x-y|}\Big)^p|f(x)-f(y)|^p \frac{dxdy}{d(x,y)^{2N}}\right)^{\frac1p}.$$
Using Theorem \ref{sobolev comparison theorem}, we conclude that
$$\|F\|_{L_p(X\times X,H^2(\mu\times\mu))}\leq c_{R,k}\|f\|_{\dot{W}^{\frac{N}{p},p}(\mathbb{R}^N)}.$$
Obviously, the integral kernel of the operator $[R_{{\rm Dunkl},j},M_f]$ is $FH.$ The assertion follows now from Lemma \ref{s4 second lemma}.
\end{proof}


\section{Proof of Theorem \ref{main theorem} \eqref{mtb}}\label{lower estimate section}


\begin{lemma}\label{nondegenerate corollary} Let $Q\in\mathcal{D}.$ We have
$$K_j(x,y)\geq \frac{c_k}{w(Q)},\quad x\in 3Q,\quad y\in Q+3s(Q)e_j.$$
\end{lemma}
\begin{proof} {\color{red} TO BE DONE. Follows from Lemma \ref{nondegenerate}.}
\end{proof}

\begin{lemma}\label{improved zhijie lemma} Let real-valued $f\in L_1^{{\rm loc}}(\mathbb{R}^N)$ be such that the commutator $[R_{{\rm Dunkl},j},M_f]$ is bounded on $L_2(\mathbb{R}^N,w(x)dx).$ For every $Q\in\mathcal{D},$ there exist $\xi_Q$ and $\eta_Q$ supported in $7Q$ such that $\|\xi_Q\|_{\infty},\|\eta_Q\|_{\infty}\leq m(Q)^{-\frac12}$ and such that
$$\frac1{w(3Q)}\int_{3Q}|f(y)-f_{3Q}|w(y)dy\leq c_k|\langle[R_{{\rm Dunkl},j},M_f]\xi_Q,\eta_Q\rangle_{L_2(\mathbb{R}^N,w(x)dx)}|.$$
\end{lemma}
\begin{proof} Let $\alpha$ be the median of $f$ on $3Q$ and $\beta$ be the median of $f$ on $Q+3s(Q)e_j.$ Assume for definiteness $\alpha\geq\beta.$ Set
$$A_1=\{x\in 3Q:\ f(x)\geq\alpha\},\quad B_1=\{x\in Q+3s(Q)e_j:\ f(x)\leq\beta\},$$
$$A_2=\{x\in 3Q:\ f(x)\leq\beta\},\quad B_2=\{x\in Q+3s(Q)e_j:\ f(x)\geq\beta\},$$
$$\xi_{Q,s}=w(Q)^{-\frac12}\chi_{A_s},\quad \eta_{Q,s}=w(Q)^{-\frac12}\chi_{B_s},\quad s=1,2.$$

We have
$$\langle[R_{{\rm Dunkl},j},M_f]\xi_{Q,s},\eta_{Q,s}\rangle_{L_2(\mathbb{R}^N,w(x)dx)}=$$
$$=\frac1{w(Q)}\int_{B_s\times A_s}K_j(x,y)(f(y)-f(x))w(x)w(y)dxdy.$$

If $x\in B_1$ and $y\in A_1,$ then
$$f(y)-f(x)=(f(y)-\alpha)+(\alpha-\beta)+(\beta-f(x))=$$
$$=|f(y)-\alpha|+|\alpha-\beta|+|f(x)-\beta|\geq |f(y)-\alpha|+|\alpha-\beta|\geq0.$$
If $x\in B_2$ and $y\in A_2,$ then
$$f(y)-f(x)=(f(y)-\beta)+(\beta-f(x))=$$
$$=-|f(y)-\beta|-|f(x)-\beta|\leq -|f(y)-\beta|\leq 0.$$

Thus,
$$\langle[R_{{\rm Dunkl},j},M_f]\xi_{Q,1},\eta_{Q,1}\rangle_{L_2(\mathbb{R}^N,w(x)dx)}\geq $$
$$\geq \frac1{w(Q)}\int_{B_1\times A_1}K_j(x,y)\big(|f(y)-\alpha|+|\alpha-\beta|\big)w(x)w(y)dxdy\geq $$
$$\geq\frac{c_k}{w^2(Q)}\int_{B_1\times A_1}\big(|f(y)-\alpha|+|\alpha-\beta|\big)w(x)w(y)dxdy=$$
$$=\frac{c_k}{w^2(Q)}\cdot\Big(w(B_1)\cdot\int_{A_1}|f(y)-\alpha|w(y)dy+|\alpha-\beta|w(B_1)w(A_1)\Big)\geq $$
$$=\frac{c_k}{w^2(Q)}\cdot\Big(\frac12w(Q)\cdot\int_{A_1}|f(y)-\alpha|w(y)dy+|\alpha-\beta|\cdot \frac14w^2(Q)\Big)\geq $$
$$\geq\frac{c_k}{4w(Q)}\cdot\Big(\int_{A_1}|f(y)-\alpha|w(y)dy+|\alpha-\beta|w(Q)\Big)$$
and
$$-\langle[R_{{\rm Dunkl},j},M_f]\xi_{Q,2},\eta_{Q,2}\rangle_{L_2(\mathbb{R}^N,w(x)dx)}\geq $$
$$\geq \frac1{w(Q)}\int_{B_2\times A_2}K_j(x,y)|f(y)-\beta|w(x)w(y)dxdy\geq $$
$$\geq\frac{c_k}{w^2(Q)}\int_{B_2\times A_2}|f(y)-\beta|w(x)w(y)dxdy=$$
$$=\frac{c_k}{w^2(Q)}\cdot w(B_2)\cdot\int_{A_2}|f(y)-\beta|w(y)dy\geq $$
$$=\frac{c_k}{2w(Q)}\int_{A_2}|f(y)-\beta|w(y)dy.$$

It remains to note that
$$\int_{3Q}|f(y)-\alpha|w(y)dy=\int_{A_1}|f(y)-\alpha|w(y)dy+$$
$$+\int_{A_2}|f(y)-\alpha|w(y)dy+\int_{3Q\backslash(A_1\cup A_2)}|f(y)-\alpha|w(y)dy\leq$$
$$\leq\int_{A_1}|f(y)-\alpha|w(y)dy+\int_{A_2}|f(y)-\beta|w(y)dy$$
$$+\int_{A_2}|\alpha-\beta|w(y)dy+\int_{3Q\backslash(A_1\cup A_2)}|\alpha-\beta|w(y)dy\leq$$
$$\leq\int_{A_1}|f(y)-\alpha|w(y)dy+\int_{A_2}|f(y)-\beta|w(y)dy+2|\alpha-\beta|w(3Q).$$
Combining the last $3$ equations, we conclude
$$\frac{c_k'}{w(Q)}\int_{3Q}|f(y)-\alpha|w(y)dy\leq\max_{s=1,2}|\langle[R_{{\rm Dunkl},j},M_f]\xi_{Q,s},\eta_{Q,s}\rangle_{L_2(\mathbb{R}^N,w(x)dx)}|.$$

If
$$|\langle[R_{{\rm Dunkl},j},M_f]\xi_{Q,1},\eta_{Q,1}\rangle_{L_2(\mathbb{R}^N,w(x)dx)}|\geq |\langle[R_{{\rm Dunkl},j},M_f]\xi_{Q,2},\eta_{Q,2}\rangle_{L_2(\mathbb{R}^N,w(x)dx)}|,$$
then we set $\xi_Q=\xi_{Q,1}$ and $\eta=\eta_{Q,1}.$ Otherwise, set $\xi_Q=\xi_{Q,2}$ and $\eta_Q=\eta_{Q,2}.$ Thus,
$$\frac{c_k'}{w(Q)}\int_{3Q}|f(y)-\alpha|w(y)dy\leq|\langle[R_{{\rm Dunkl},j},M_f]\xi_Q,\eta_Q\rangle_{L_2(\mathbb{R}^N,w(x)dx)}|.$$
Hence,
$$c_k'|f_{3Q}-\alpha|\leq|\langle[R_{{\rm Dunkl},j},M_f]\xi_Q,\eta_Q\rangle_{L_2(\mathbb{R}^N,w(x)dx)}|$$
and, finally,
$$\frac{c_k'}{w(Q)}\int_{3Q}|f(y)-f_{3Q}|w(y)dy\leq 2|\langle[R_{{\rm Dunkl},j},M_f]\xi_Q,\eta_Q\rangle_{L_2(\mathbb{R}^N,w(x)dx)}|.$$
\end{proof}

\begin{lemma}\label{standard} Suppose that $1<p<\infty$ and $f\in L_1^{{\rm loc}}(\mathbb{R}^N).$ If $[R_{{\rm Dunkl},j},M_f]\in\mathcal{L}_p(L_2(\mathbb{R}^N,w(x)dx))$ for some $j\in\{1,2,\cdots,N\},$ then
$$\Big\|\Big\{\frac1{w(3Q)}\int_{3Q}|f(x)-f_{3Q}|w(x)dx\Big\}_{Q\in\mathcal{D}}\Big\|_{l_p(\mathcal{D})}\leq$$
$$\leq c_k\|[R_{{\rm Dunkl},j},M_f]\|_{\mathcal{L}_p(L_2(\mathbb{R}^N,w(x)dx))}.$$
If $[R_{{\rm Dunkl},j},M_f]\in\mathcal{L}_{N,\infty}(L_2(\mathbb{R}^N,w(x)dx))$ for some $j\in\{1,2,\cdots,N\},$ then
$$\Big\|\Big\{\frac1{w(3Q)}\int_{3Q}|f(x)-f_{3Q}|w(x)dx\Big\}_{Q\in\mathcal{D}}\Big\|_{l_{N,\infty}(\mathcal{D})}\leq$$
$$\leq c_k\|[R_{{\rm Dunkl},j},M_f]\|_{\mathcal{L}_{N,\infty}(L_2(\mathbb{R}^N,w(x)dx))}.$$
\end{lemma}
\begin{proof} We prove only the first assertion. The proof of the second one follows {\it mutatis mutandi}.

Let $\{\xi_Q\}_{Q\in\mathcal{D}}$ and $\{\eta_Q\}_{Q\in\mathcal{D}}$ be functions constructed in Lemma \ref{improved zhijie lemma}. By Lemma \ref{NWO lower estimate},
$$\Big\{\langle [R_{{\rm Dunkl},j},M_f]\xi_Q,\eta_Q\rangle_{L_2(\mathbb{R}^N,w(x)dx)}\Big\}_{Q\in\mathcal{D}}\in l_p.$$
The first assertion follows now from Lemma \ref{improved zhijie lemma}.
\end{proof}

\begin{lemma}\label{for reverse holder first lemma} Let $Q$ be a cube in $\mathbb{R}^N$ (with sides parallel to coordinate hyperplanes). We have
$$\sup_{x\in Q}w(x)\leq \frac{c_k}{m(Q)}\int_Qw(x)dx.$$
\end{lemma}
\begin{proof} If $x=c_Q+y,$ $y\in [-\frac12s(Q),\frac12s(Q)]^d,$ then
$$|\langle x,v\rangle|\leq |\langle c_Q,v\rangle|+r|\langle y,v\rangle|\leq |\langle c_Q,v\rangle|+\frac{\sqrt{d}}{2}s(Q).$$
Thus,
$$w(x)=\prod_{v\in R}|\langle x,v\rangle|^{k(v)}\leq\prod_{v\in R}(|\langle c_Q,v\rangle|+\frac{\sqrt{d}}{2}s(Q))^{k(v)}.$$
In particular,
$$LHS\leq \prod_{v\in R}(|\langle c_Q,v\rangle|+\frac{\sqrt{d}}{2}s(Q))^{k(v)}.$$

Fix chamber $\Gamma$ such that $c_Q\in\Gamma.$ We have
$$\int_Qw(x)dx\geq\int_{c_Q+(B(0,\frac12s(Q))\cap\Gamma)}w(x)dx=\int_{B(0,\frac12s(Q))\cap\Gamma}w(c_Q+y)dy.$$
If $y\in\Gamma,$ then
$$|\langle c_Q+y,v\rangle|=|\langle c_Q,v\rangle|+|\langle y,v\rangle|,\quad v\in R.$$
Thus,
$$w(c_Q+y)=\prod_{v\in R}(|\langle c_Q,v\rangle|+|\langle y,v\rangle|)^{k(v)}.$$
Hence,
$$\int_Qw(x)dx\geq\int_{B(0,\frac12s(Q))\cap\Gamma}\prod_{v\in R}(|\langle c_Q,v\rangle|+|\langle y,v\rangle|)^{k(v)}dy.$$
Fix a closed cone $\Gamma'\subset\Gamma$ such that (i) interior of $\Gamma'$ is non-empty and (ii) $\partial\Gamma\cap\partial\Gamma'=\{0\}.$ We have
$$|\langle y,v\rangle|\geq c_{\Gamma'}|y|,\quad y\in\Gamma',\quad v\in R.$$
We, therefore, have
$$\int_Qw(x)dx\geq\int_{B(0,\frac12s(Q))\cap\Gamma'}\prod_{v\in R}(|\langle c_Q,v\rangle|+c_{\Gamma'}|y|)^{k(v)}dy\geq$$
$$\geq \int_{\substack{y\in\Gamma'\\ \frac14s(Q)\leq|y|\leq \frac14s(Q)}}\prod_{v\in R}(|\langle c_Q,v\rangle|+\frac12c_{\Gamma'}s(Q))^{k(v)}dy=$$
$$=s(Q)^N\prod_{v\in R}(|\langle c_Q,v\rangle|+\frac14c_{\Gamma'}s(Q))^{k(v)}\times (2^{-N}-4^{-N})\int_{\substack{y\in\Gamma'\\ |y|\leq1}}dy.$$
Combining this with the preceding paragraph, we complete the proof.
\end{proof}

\begin{lemma}\label{for reverse holder second lemma} Suppose $1<t<1+\frac1{|R|\cdot\max_{v\in R}k(v)}.$ Let $Q$ be a cube in $\mathbb{R}^N$ (with sides parallel to coordinate hyperplanes). We have
$$\Big(\frac{\int_Qw^{1-t}(x)dx}{\int_Qw(x)dx}\Big)^{\frac1t}\leq c_{t,k} \frac{m(Q)}{\int_Qw(x)dx}.$$
\end{lemma}
\begin{proof} If $v\in R$ is such that
$$\sqrt{d}s(Q)\leq |\langle c_Q,v\rangle|,$$
then
$$|\langle x,v\rangle|\geq |\langle c_Q,v\rangle|-|\langle x-c_Q,v\rangle|\geq |\langle c_Q,v\rangle|-\frac{\sqrt{d}}{2}s(Q)\geq \frac12|\langle c_Q,v\rangle|,\quad x\in Q.$$
In this case,
$$\int_Q|\langle x,v\rangle|^{(1-t)|R|k(v)}dx\leq 2^{(t-1)|R|k(v)}|\langle c_Q,v\rangle|^{(1-t)|R|k(v)}\cdot\int_Qdx=$$
$$=2^{(t-1)|R|k(v)}|\langle x_0,v\rangle|^{(1-t)|R|k(v)}\cdot s(Q)^N.$$

If $v\in R$ is such that
$$|\langle c_Q,v\rangle|\leq \sqrt{d}s(Q),$$
then let $y_0=\langle c_Q,v\rangle$ and $z_0=c_Q-y_0v.$ Let us introduce new variables $y=\langle x,v\rangle$ and $z=x-yv$ (it varies in the plane orthogonal to $v$). We write
$$\int_Q|\langle x,v\rangle|^{(1-t)|R|k(v)}dx\leq\int_{|x-c_Q|\leq\frac{\sqrt{d}}{2}s(Q)}|\langle x,v\rangle|^{(1-t)|R|k(v)}dx\leq$$
$$\leq \int_{|y-y_0|,|z-z_0|\leq \frac{\sqrt{d}}{2}r}|y|^{(1-t)|R|k(v)}dydz=$$
$$=(\frac{\sqrt{d}}{2}s(Q))^{N-1}{\rm Vol}(\mathbb{B}^{N-1})\cdot\int_{|y-y_0|\leq \frac{\sqrt{d}}{2}s(Q)}|y|^{(1-t)|R|k(v)}dy\leq$$
$$=(\frac{\sqrt{d}}{2}s(Q))^{N-1}{\rm Vol}(\mathbb{B}^{N-1})\cdot\int_{|y|\leq \frac{3\sqrt{d}}{2}s(Q)}|y|^{(1-t)|R|k(v)}dy=$$
$$=s(Q)^{N+(1-t)|R|k(v)}\cdot (\frac{\sqrt{d}}{2})^{N-1}{\rm Vol}(\mathbb{B}^{N-1})\cdot\int_{|y|\leq \frac{3\sqrt{d}}{2}}|y|^{(1-t)|R|k(v)}dy.$$

Combining the estimates in the preceding paragraphs, we write
$$\int_Q|\langle x,v\rangle|^{(1-t)|R|k(v)}dx\leq c_1s(Q)^N\big(\max\{|\langle c_Q,v\rangle|,r\}\big)^{(1-t)|R|k(v)}.$$
By H\"older inequality,
$$\int_Qw^{1-t}(x)dx\leq \prod_{v\in R}\Big(\int_Q|\langle x,v\rangle|^{(1-t)|R|k(v)}dx\Big)^{\frac1{|R|}}\leq$$
$$\leq c_2s(Q)^N\prod_{v\in R}\big(\max\{|\langle c_Q,v\rangle|,s(Q)\}\big)^{(1-t)k(v)}.$$
The assertion follows easily.
\end{proof}

\begin{lemma}\label{hytonen verification lemma} Measures $dx$ and $w(x)dx$ (and distance $(x,y)\to |x-y|_{\infty}$) satisfy the conditions in Proposition 1.2 in \cite{Hytonen-new}.
\end{lemma}
\begin{proof} That both measures are doubling is obvious.

Let $Q$ be a cube in $\mathbb{R}^N$ (with sides parallel to coordinate hyperplanes). By Lemmas \ref{for reverse holder first lemma} and \ref{for reverse holder second lemma},
$$\Big(\frac1{m(Q)}\int_Qw^t(x)dx\Big)^{\frac1t}\leq\frac{c_k}{m(Q)}\int_{\Gamma}w(x)dx,\quad t>0.$$
$$\Big(\frac{\int_Qw^{1-t}(x)dx}{\int_Qw(x)dx}\Big)^{\frac1t}\leq c_G \frac{m(Q)}{\int_Qw(x)dx},\quad 1<t<1+\frac1{|R|\cdot\max_{v\in R}k(v)}.$$
The assertion follows now from Proposition 3.6 in \cite{Hytonen-new}.
\end{proof}

\begin{proof}[Proof of Theorem \ref{main theorem} \eqref{mtb}] The left hand side in Lemma \ref{standard} is the oscillation semi-norm of the function $f$ with respect to measure $w(x)dx$ and distance $(x,y)\to |x-y|_{\infty}.$ So, the oscillation semi-norm of the function $f$ with respect to measure $w(x)dx$ and distance $(x,y)\to |x-y|_{\infty}$ is controlled by the $\|\cdot\|_{\mathcal{L}_p(L_2(\mathbb{R}^N,w(x)dx))}$-norm of the commutator $[R_{{\rm Dunkl},j},M_f].$

Lemma \ref{hytonen verification lemma} asserts that measures $dx$ and $w(x)dx$ (and distance $(x,y)\to |x-y|_{\infty}$) satisfy the conditions in Proposition 1.2 in \cite{Hytonen-new}. Appealing to Proposition 1.2 in \cite{Hytonen-new}, we conclude that the oscillation semi-norm of the function $f$ with respect to measure $dx$ and distance $(x,y)\to |x-y|_{\infty}$ is controlled by the one with respect to measure $w(x)dx$ and distance $(x,y)\to |x-y|_{\infty}.$ By the preceding paragraph, the oscillation semi-norm of the function $f$ with respect to measure $dx$ and distance $(x,y)\to |x-y|_{\infty}$ is controlled by the $\|\cdot\|_{\mathcal{L}_p(L_2(\mathbb{R}^N,w(x)dx))}$-norm of the commutator $[R_{{\rm Dunkl},j},M_f].$

That oscillation semi-norm of the function $f$ with respect to measure $dx$ and distance $(x,y)\to |x-y|_{\infty}$ is equivalent to the Sobolev semi-norm of the function $f$ is a standard fact. The proof (in much more general setting) is available in Proposition 1.29 in \cite{HKarxiv}. The assertion follows now from the preceding paragraph. 
\end{proof}

\begin{proof}[Proof of Theorem \ref{main theorem} \eqref{mtd}] The argument follows the one in the proof of Theorem \ref{main theorem} \eqref{mtb} {\it mutatis mutandi}.
\end{proof}


\section{Constant characterization}
\setcounter{equation}{0}

\begin{theorem} If $f\in L_1^{{\rm loc}}(\mathbb{R}^N,w(x)dx)$ is such that $$[R_{{\rm Dunkl},j},M_f]\in\mathcal{L}_p(L_2(\mathbb{R}^N,w(x)dx))$$
for some $p\leq d,$ then $f$ is constant.
\end{theorem}
\begin{proof} By Lemma \ref{hytonen verification lemma}, measures $dx$ and $w(x)dx$ and operator $T_{\Gamma,j}=M_{\chi_{\Gamma}}R_{{\rm Dunkl},j}M_{\chi_{\Gamma}}$ satisfy {\color{red} on $\Gamma$} the assumptions in Theorem 1.4 in \cite{Hytonen-new}. By assumption, $[T_{\Gamma,j},M_f]\in\mathcal{L}_p(L_2(\Gamma,w(x)dx)).$ By Theorem 1.4 in \cite{Hytonen-new}, $f$ is constant on $\Gamma.$

Hence, $f$ is constant on every chamber. Let $f=c_1$ on $\Gamma_1$ and $f=c_2$ on $\Gamma_2.$ On the other hand, $[R_{{\rm Dunkl},j},M_f]\in\mathcal{L}_{2N}(L_2(\mathbb{R}^N,w(x)dx)).$ By Theorem \ref{main theorem} \eqref{mtb}, $f\in\dot{W}^{\frac12,2N}(\mathbb{R}^N).$  We have
$$\int_{\Gamma_1\times\Gamma_2}\frac{|c_1-c_2|^{2N}dxdy}{|x-y|^{2N}}=\int_{\Gamma_1\times\Gamma_2}\frac{|f(x)-f(y)|^{2N}dxdy}{|x-y|^{2N}}\leq$$
$$\leq\int_{\mathbb{R}^N\times\mathbb{R}^N}\frac{|f(x)-f(y)|^{2N}dxdy}{|x-y|^{2N}}=\|f\|_{\dot{W}^{\frac12,2N}(\mathbb{R}^N)}<\infty.$$
It follows that $c_1=c_2.$ Since chambers $\Gamma_1$ and $\Gamma_2$ are arbitrary, the assertion follows.
\end{proof}


 \noindent
 {\bf Acknowledgements:}

%The authors would like to thank Rambau J\"{o}rg, Christian Kipp, Jes\'{u}s De Loera, Santos Leal  Francisco, Arnau Padrol for helpful discussion on triangulation of polyhedral cone.
Z. Fan is supported by the Guangdong Basic and Applied Basic Research Foundation (No. 2023A1515110879). J. Li is supported by the Australian Research Council (ARC) through the research grant DP220100285.

\bibliographystyle{plain}
\bibliography{references}

\end{document}
